\documentclass[../tesis.tex]{subfiles}
\begin{document}

\chapter*{Glosario}

\begin{description}[font=\bfseries, leftmargin=4cm, style=nextline]

    \item[Absorciometría de Rayos X de Energía Dual (DXA)]
    Técnica de imagen de baja dosis radiactiva que mide la Densidad Mineral Ósea (DMO) y la
    composición corporal a partir de la atenuación diferencial de dos haces de rayos X con
    distinta energía. Es el estándar de oro para el diagnóstico de osteoporosis según los
    criterios de la OMS (1994).

    \item[Análisis Estructural de Cadera (HSA)]
    Conjunto de métricas geométricas derivadas del escaneo DXA que caracterizan la resistencia
    biomecánica del cuello femoral: Área de Sección Transversal (CSA), Índice de Sección de
    Módulo Compuesto (CSMI), Longitud del Eje de la Cadera (HAL) y ancho cortical. Permite
    estimar la fragilidad estructural de forma independiente a la DMO puntual.

    \item[Densidad Mineral Ósea (DMO / BMD)]
    Masa de mineral óseo por unidad de área proyectada (g/cm²), estimada por DXA en sitios
    esqueléticos de interés clínico: cuello femoral, cadera total y columna lumbar (L1--L4).
    Es la variable primaria sobre la que se calculan el T-score y el Z-score diagnósticos.

    \item[Evaluación de Fractura Vertebral (VFA)]
    Modalidad del DXA que genera una imagen lateral de la columna dorso-lumbar (T6--L4) con
    dosis de radiación mínima, permitiendo la detección semicuantitativa de deformidades
    vertebrales y micro-fracturas silenciosas mediante la clasificación de Genant (1993).

    \item[Fractura por fragilidad]
    Fractura producida por un mecanismo de baja energía, equivalente a una caída desde la propia
    altura o por esfuerzo mínimo. Es indicativa de fragilidad esquelética subyacente y constituye
    el desenlace clínico primario que los modelos de riesgo buscan predecir.

    \item[FRAX]
    Herramienta de cálculo probabilístico desarrollada por la Universidad de Sheffield (Kanis
    et~al.) que combina la DMO del cuello femoral con once factores clínicos para estimar la
    probabilidad absoluta de fractura mayor osteoporótica y de fractura de cadera a 10 años.
    Se calibra por país a partir de datos de epidemiología de fracturas y mortalidad locales.

    \item[Modelo multimodal]
    Modelo predictivo que integra simultáneamente variables de distinta naturaleza o fuente
    (p.ej. DMO + HSA + VFA + composición corporal + datos clínicos tabulares), en lugar de
    basarse en una única señal. La fusión de modalidades permite capturar interacciones que
    ningún marcador individual puede reflejar por sí solo.

    \item[Osteosarcopenia]
    Condición musculoesquelética sinérgica definida por la coexistencia de baja masa ósea
    (osteopenia/osteoporosis) y sarcopenia. Implica un riesgo de caídas y fracturas
    amplificado respecto al de cada condición por separado, y puede estar presente aun cuando
    el T-score aislado se sitúe en rango normal.

    \item[\textit{Pipeline} metodológico]
    Secuencia estructurada de etapas de procesamiento de datos y modelado, desde la ingesta
    y curación de la fuente de datos hasta la generación de un reporte interpretable para el
    clínico.     En esta investigación: ETL $\rightarrow$ Integración multimodal $\rightarrow$
    Modelo transversal $\rightarrow$ Interpretabilidad $\rightarrow$ Validación.

    \item[SHAP (\textit{SHapley Additive exPlanations})]
    Marco de interpretabilidad basado en la teoría de juegos cooperativos (valores de Shapley)
    que asigna a cada variable predictora una contribución marginal cuantificada a la predicción
    individual del modelo. Garantiza tres propiedades axiomáticas: precisión local (la suma de
    contribuciones individuales iguala la predicción del modelo), \textit{missingness} (variables
    ausentes reciben contribución cero) y consistencia (si una variable contribuye más bajo
    cualquier coalición posible, su valor SHAP no disminuye).

    \item[T-score]
    Número de desviaciones estándar en las que la DMO de un paciente se aleja de la media de una
    población adulta joven y sana de referencia, medida en el mismo sitio esquelético con el mismo
    tipo de equipo. Criterio diagnóstico OMS: $\leq -2{,}5$ osteoporosis; entre $-2{,}5$ y
    $-1{,}0$ osteopenia; $\geq -1{,}0$ normal.

    \item[Z-score]
    Número de desviaciones estándar en las que la DMO de un paciente se aleja de la media de
    individuos de su misma edad, sexo y grupo étnico de referencia. A diferencia del T-score,
    contextualiza la pérdida ósea respecto al envejecimiento esperado y es de utilidad clínica
    preferente en pacientes premenopáusicas y adultos menores de 50 años.

\end{description}
    \newpage

\end{document}
